%%
%% CSD Assignment 1 - Final Technical Report
%% Based on the ACM "acmart" manuscript (single-column) style.
%%
\documentclass[manuscript,screen,review,nonacm]{acmart}
\usepackage{tikz}
\usetikzlibrary{shapes.geometric, arrows.meta, positioning, fit, backgrounds, decorations.pathreplacing}

\AtBeginDocument{%
  \providecommand\BibTeX{{Bib\TeX}}}

\begin{document}

%% =================== TITLE & AUTHOR BLOCK =========================
\title{AI-based Village Pond Planning System}
\subtitle{CSD Assignment 1 -- Final Technical Report}

\author{Harshita Patidar -- 12340920}
\affiliation{%
  \institution{Indian Institute of Technology, Bhilai}%
  \department{Department of Computer Science and Engineering}%
  \country{India}}
\email{harshitap@iitbhilai.ac.in}

\renewcommand{\shortauthors}{Patidar}

%% =================== ABSTRACT ======================================
\begin{abstract}
Rural water scarcity in India is a persistent challenge, particularly in
rain-fed agricultural regions where ponds remain the most cost-effective
form of water storage. Manual identification of optimal pond sites is
labour-intensive, inconsistent, and often ignores quantitative terrain and
hydrological evidence. This report describes a web-based Village Pond
Planning System that automates site selection using geospatial terrain
analysis, satellite rainfall data, and classical civil engineering design
standards. The system accepts either a user-supplied KMZ contour map or a
user-drawn map rectangle (using free SRTM satellite elevation tiles), and
runs a stepwise pipeline: DEM interpolation, Priority-Flood sink filling,
D8 flow routing, OSM-based land exclusion, candidate depression scoring,
SCS-CN runoff estimation, and IS~5477-compliant pond dimensioning. Results
are visualised on an interactive Leaflet.js map with a full engineering
sidebar. On a representative test site in central India (catchment area
19.5~ha, annual rainfall 1,458~mm), the system recommended a pond with
1.84~m water depth, 61,400~m$^3$ live storage, and a wet-season fill
probability of 100\%. End-to-end analysis completes in under 90~seconds.
\end{abstract}

\keywords{Village Pond Planning, Geospatial Analysis, Rainwater
  Harvesting, Catchment Estimation, SCS-CN, Web GIS, REST API, DEM,
  IS~5477}

\maketitle

%% =====================================================================
\section{Introduction}
\label{sec:intro}

India loses an estimated 70\% of its annual rainfall as runoff due to
the absence of local water retention infrastructure. Village ponds,
functioning as distributed rainwater harvesting tanks, are among the most
effective interventions: they recharge groundwater, support irrigated
agriculture in the dry season, and provide livestock and domestic water
supplies. Despite proven benefits, new pond construction is impeded by the
difficulty of site selection --- a process that currently requires
expensive surveying, manual terrain reading, and domain expertise rarely
available at the village level.

This project develops a \emph{web-based Village Pond Planning System} that
automates the entire site-selection workflow using freely available
geospatial data, satellite rainfall records, and classical hydrological
methods. A village administrator uploads a topographic contour map (or
simply draws a bounding box on the map), and the system returns ranked pond
candidates, a delineated catchment watershed, historical rainfall
statistics, estimated runoff volumes, and a pond design that complies with
IS~5477 and IS~12169 standards --- all within 90~seconds.

The rest of this report describes the system requirements
(Section~\ref{sec:requirements}), the architecture
(Section~\ref{sec:architecture}), the algorithms
(Section~\ref{sec:methodology}), implementation details
(Section~\ref{sec:implementation}), evaluation with real outputs
(Section~\ref{sec:results}), limitations (Section~\ref{sec:discussion}),
and an AI usage declaration (Section~\ref{sec:ai-usage}).

\subsection{Motivation}

Site selection by hand is hard for three compounding reasons. First,
natural drainage basins rarely coincide with administrative boundaries, so
a pond that looks central on a revenue map may capture only a fraction of
the upstream rainfall. Second, rainfall data is fragmented across IMD
district records, NASA POWER, and Open-Meteo; aggregating a coherent
10-year series for an arbitrary coordinate requires API knowledge that
most field officials lack. Third, candidate sites must be cross-checked
against existing water bodies, buildings, and land-tenure records to avoid
legal and engineering conflicts. Automating these steps with reproducible,
open-source tools directly reduces the expertise barrier for rural water
managers.

\subsection{Scope of the Project}

The system covers:
\begin{itemize}
  \item Ingestion of user-supplied KMZ contour maps \emph{or} free SRTM
        tiles fetched by bounding-box selection on the map;
  \item 2D DEM construction at $\approx$2~m resolution (KMZ path) or
        $\approx$38~m (SRTM tile path);
  \item Automated depression detection, D8 flow routing, and ranked
        candidate site generation;
  \item OSM-based exclusion of water bodies and built-up land;
  \item 10-year monthly rainfall retrieval (Open-Meteo / NASA POWER
        fallback);
  \item SCS-CN runoff estimation and IS~5477 pond dimensioning;
  \item Interactive map visualisation and a structured JSON API.
\end{itemize}

The system does \emph{not} cover structural embankment engineering,
real-time flood forecasting, land-ownership verification, or construction
cost estimation.

%% =====================================================================
\section{Problem Statement and Requirements}
\label{sec:requirements}

\begin{table}[h]
  \caption{Functional requirements and their implementation}
  \label{tab:requirements}
  \begin{tabular}{p{3.6cm}p{3.6cm}}
    \toprule
    Requirement & Implemented in \\
    \midrule
    Satellite imagery display & \texttt{src/api/imagery.py} + Leaflet basemap switcher (3 providers) \\
    Contour map / terrain input & \texttt{src/terrain/kml\_source.py} (KMZ) + \texttt{src/dem/from\_tiles.py} (SRTM) \\
    Available-land identification & \texttt{src/catchment/water\_exclusion.py} + \texttt{land\_exclusion.py} (OSM Overpass) \\
    Catchment area estimation & \texttt{src/hydrology/watershed.py} (D8 BFS delineation) \\
    Historical rainfall query & \texttt{src/catchment/rainfall\_service.py} (Open-Meteo; NASA POWER fallback) \\
    Runoff volume estimation & \texttt{src/hydrology/runoff.py} (SCS-CN method) \\
    Pond depth / storage & \texttt{src/catchment/pond\_design.py} (IS~5477) \\
    Combined overlay / results & \texttt{frontend/index.html} (Leaflet + Chart.js) \\
    \bottomrule
  \end{tabular}
\end{table}

\subsection{Non-Functional Requirements}

\textbf{Performance.} End-to-end analysis must complete in under 120~s.
Rainfall API calls must complete within 30~s with automatic retry. SRTM
tile download for a 10~km$^2$ area must complete within 10~s.

\textbf{Reliability.} Every external API call (OSM Overpass, Open-Meteo,
NASA POWER) is wrapped in a fail-open handler: if the call fails after
retries, the system returns a degraded-but-valid result with
machine-readable warning codes rather than crashing.

\textbf{Scalability.} CPU-intensive computation runs in a
\texttt{run\_in\_threadpool} worker to keep the async event loop free for
concurrent HTTP requests. Horizontal scaling via multiple uvicorn workers
is supported by the stateless design.

\textbf{Usability.} The UI must operate in a modern desktop browser with
no plugins. Area selection must prevent analyses larger than 10~km$^2$
to protect server resources (configurable via \texttt{.env}).

%% =====================================================================
\section{System Architecture and High-Level Design}
\label{sec:architecture}

The system follows a three-tier architecture: a Vanilla JS single-page
application (SPA) frontend, a Python/FastAPI REST backend, and optional
PostgreSQL persistence. Two analysis entry points
(\texttt{POST /analyzeContour} and \texttt{POST /analyzeArea}) converge
at a shared pipeline from sink-filling onward, avoiding code duplication.
The only structural difference is how the DEM is built --- from IDW
interpolation of KMZ contours, or from stitching Terrarium RGB PNG tiles.



% \caption{System architecture. Both input paths (KMZ upload and Draw Area
%   SRTM tiles) produce a DEM that enters the shared hydrology pipeline.
%   Dashed purple arrows denote external API calls, all wrapped in fail-open
%   handlers with retry. The result is returned to the browser and
%   optionally persisted in PostgreSQL.}
% \label{fig:architecture}
% \Description{Vertical flow diagram: Browser at top splits to KMZ Upload
%   and Draw Area paths, each producing a DEM (IDW 2 m or SRTM 38 m).
%   Both merge into FastAPI, then a shared pipeline of Priority-Flood,
%   D8 flow routing, OSM exclusion, candidate scoring, rainfall retrieval,
%   SCS-CN runoff, IS-5477 pond design, and BFS watershed delineation,
%   outputting AnalysisResult JSON back to the browser and optionally
%   to PostgreSQL. External APIs shown on the right.}

 \begin{figure}[h]
   \centering
    \includegraphics[width=0.5\linewidth]{Screenshot from 2026-09-27 22-22-58.png}
   \caption{High-level system architecture. Both analysis paths produce a
     DEM object consumed by the shared pipeline.}
   \label{fig:architecture}
   \Description{Flow diagram: browser SPA sends POST /analyzeContour or
     POST /analyzeArea to FastAPI. KMZ path: KMLTerrainSource + IDW to DEM.
     Tile path: Terrarium tile download + decode to DEM. Shared pipeline:
     Priority-Flood, D8 routing, OSM masks, candidates, rainfall, SCS-CN,
    IS-5477, BFS watershed, polygonise, AnalysisResult JSON to browser.}
 \end{figure}


\subsection{Technology Stack}

\begin{table}[h]
  \caption{Technology stack}
  \label{tab:stack}
  \begin{tabular}{p{2.5cm}p{4.7cm}}
    \toprule
    Layer & Technology \\
    \midrule
    Language & Python 3.12 \\
    Web framework & FastAPI 0.111 + uvicorn 0.29 \\
    Geospatial & NumPy, SciPy, PyProj, Shapely, pysheds, Rasterio \\
    KML parsing & fastkml 0.12, lxml 5.2 \\
    Tile DEM & Pillow 10.3, httpx 0.27 \\
    Rainfall & Open-Meteo ERA5-Land (free); NASA POWER MERRA-2 (free) \\
    Terrain tiles & AWS Terrarium (SRTM, public domain, no key) \\
    OSM exclusion & Overpass API (free, no key) \\
    Persistence & SQLModel + asyncpg + PostgreSQL (optional) \\
    Config & pydantic-settings, \texttt{.env} file \\
    Frontend & HTML/CSS/JS, Leaflet.js, Chart.js \\
    CI & GitHub Actions, black, ruff, pre-commit \\
    \bottomrule
  \end{tabular}
\end{table}

No paid APIs, proprietary datasets, or ML models are used. All data
sources are freely available and reproducible.

%% =====================================================================
\section{Methodology}
\label{sec:methodology}

\subsection{Terrain and Elevation Analysis}

\textbf{KMZ path.} KML/KMZ contour files are parsed with fastkml + lxml
to extract elevation \texttt{LineString} rings. Each ring is reprojected
from WGS84 to the local UTM zone (auto-detected from centroid longitude).
Contour lines are densified to 1~m vertex spacing to form a dense 3D
point cloud. A 2D regular grid (default 2~m $\times$ 2~m) is populated
by \emph{Inverse Distance Weighting} (IDW) using a SciPy \texttt{KDTree}:
each grid cell receives a $1/d^2$-weighted average elevation from its $k$
nearest contour points. Edge cells are filled by nearest-neighbour.

\textbf{SRTM tile path.} The user draws a bounding box on the map.
\texttt{src/dem/from\_tiles.py} downloads Terrarium-format PNG tiles from
the AWS Open Data programme at zoom level 12 ($\approx$38~m/pixel).
Elevation is decoded as $e = (R\times256 + G + B/256) - 32768$. Tiles
are stitched, clipped to the bounding box, and reprojected to UTM,
yielding a DEM object structurally identical to the KMZ output.

In both paths the DEM passes through the \emph{Planchon-Darboux
Priority-Flood} algorithm (pysheds) to fill artificial pits, producing a
hydrologically conditioned filled DEM. Slope is computed from the filled
DEM.

\subsection{Catchment Area Delineation}

The filled DEM is processed by pysheds' D8 flow-direction algorithm,
assigning each raster cell a pointer to its steepest downhill neighbour.
A topological traversal produces a flow-accumulation raster.

Depression bowls are identified as connected components where
$\Delta z = \text{filled\_dem} - \text{raw\_dem} > \delta_{min}$
(default 0.1~m). For each bowl, the lowest-elevation sink cell is the
pour point. Candidates failing hard veto rules (OSM water/land mask
overlap, slope $>20°$, area $< 500$~m$^2$) are discarded. Remaining
candidates are scored by normalised catchment size:
\[
  s_i = \frac{\text{flow\_accumulation}(\text{pour\_point}_i)}{\text{total DEM cells}}
\]

Watershed delineation for the selected site uses a BFS traversal of the
reversed D8 grid, collecting all upstream cells. The boolean raster is
polygonised, Douglas-Peucker simplified, and reprojected to WGS84 GeoJSON.

\subsection{Rainfall Data Integration}

After candidate selection, the system fetches 10-year monthly rainfall
totals for the selected site. The primary source is Open-Meteo ERA5-Land
(free, $\approx$9~km resolution). On failure after 3 retries, NASA POWER
MERRA-2 ($\approx$50~km) is used. The current calendar year is excluded
to ensure complete annual records. Monthly totals are aggregated into
wet-season (June--September) and dry-season statistics. Failure of both
APIs is fail-open: fields are returned as \texttt{null} with warning codes.

%% =====================================================================
\section{Implementation}
\label{sec:implementation}

\subsection{Backend and API Design}

\begin{table}[h]
  \caption{REST API endpoints}
  \label{tab:api}
  \begin{tabular}{p{3.7cm}p{3.5cm}}
    \toprule
    Endpoint & Purpose \\
    \midrule
    \texttt{POST /analyzeContour} & KMZ/KML upload $\to$ full analysis \\
    \texttt{POST /analyzeArea} & Bbox $\to$ SRTM DEM $\to$ analysis \\
    \texttt{GET /health} & Liveness probe \\
    \texttt{GET /api/imagery} & Basemap tile URL metadata \\
    \texttt{GET /api/village/search?q=} & Village registry search \\
    \texttt{GET /api/village/\{id\}/kml} & Download registered village KML \\
    \texttt{GET /api/results/\{id\}} & Retrieve saved result (requires DB) \\
    \bottomrule
  \end{tabular}
\end{table}

No authentication is implemented (demonstration use). All endpoints return
structured JSON. CPU-bound analysis is offloaded via
\texttt{run\_in\_threadpool} to keep the async event loop free. In-memory
TTL caches (24~h) are applied to OSM Overpass and rainfall responses,
reducing repeat API calls within a session to 0~ms.

\subsection{Frontend and Visualization}

The frontend is a single HTML file (\texttt{frontend/index.html}) using
Leaflet.js for the interactive map and Chart.js for the monthly rainfall
histogram. No build step or bundler is required.

\begin{figure}[h]
  \centering
   \includegraphics[width=\linewidth]{Screenshot from 2026-09-27 22-32-53.png}
  \caption{Application interface: satellite basemap with blue catchment
    polygon, blue candidate dots, red selected-site star, and dashed blue
    analysis area rectangle. The sidebar shows all results cards and the
    monthly rainfall bar chart.}
  \label{fig:ui}
  \Description{Screenshot of the Village Pond Planning System. Left sidebar
    shows four result cards (Selected Site, Catchment Area, Rainfall and
    Runoff, Recommended Pond Design) and a monthly rainfall bar chart. The
    map shows the catchment polygon and candidate markers.}
\end{figure}

Two analysis workflows are available:
\begin{enumerate}
  \item \textbf{KMZ Upload}: drag-and-drop a contour file. Contours are
        previewed on the map while the API processes.
  \item \textbf{Draw Area}: click ``Draw Area on Map'', drag a rectangle
        ($\le$10~km$^2$), click Analyze. The bounding box persists after
        results as a dashed blue reference overlay; the legend reflects
        this with a ``Selected area'' entry.
\end{enumerate}

A pour-point override button allows clicking any map location to
re-run the analysis snapped to the nearest detected candidate bowl.
A companion page (\texttt{frontend/help.html}) documents all features,
result field definitions, and data source accuracy for end users.

\subsection{Database and Storage}

Result persistence is optional and disabled by default. When
\texttt{DATABASE\_URL} is set, results are written to an
\texttt{analysis\_runs} table (PostgreSQL JSONB) \emph{after} the HTTP
response is assembled, so a DB outage never increases latency.
A \texttt{Village} registry table stores pre-registered village metadata
and KML paths for the search-and-load workflow.

%% =====================================================================
\section{CSD Themes and Topics Applied in the Project}
\label{sec:csd-themes}

\begin{table*}[t]
  \caption{Mapping of core CSD themes/topics to their use in this project}
  \label{tab:csd-themes}
  \begin{tabular}{p{2.8cm}p{5.2cm}p{6.5cm}}
    \toprule
    CSD Theme/Topic & Where used & Justification / design rationale \\
    \midrule
    API Design (REST) & \texttt{POST /analyzeContour}, \texttt{POST /analyzeArea}, \texttt{GET /api/imagery} & REST chosen for statelessness and browser compatibility. Both analysis endpoints return identical \texttt{AnalysisResult} JSON, so the frontend needs zero branching logic to display results from either path. \\
    Caching & In-memory TTL dict (24~h) for OSM XML and rainfall in \texttt{water\_exclusion.py} and \texttt{rainfall\_service.py} & Rainfall and water-feature data are static within a planning session; caching eliminates redundant API calls and reduces latency from $\approx$500~ms to 0~ms on cache hits. \\
    Concurrency / Async & \texttt{run\_in\_threadpool} wraps CPU-bound analysis; FastAPI async routes handle HTTP I/O & The pipeline is CPU-bound (IDW, pysheds BFS); offloading to a thread pool keeps the event loop free for concurrent connections. \\
    Microservices vs.\ Monolith & Modular monolith: single FastAPI app with clean internal module boundaries & For a demonstration system the operational overhead of microservices outweighs benefits. All modules are independently unit-tested; extracting a service requires only a thin HTTP wrapper. \\
    Design Patterns & Strategy (\texttt{TerrainSource} ABC) for input formats; Singleton for \texttt{Settings} and \texttt{AnalysisService}; Fail-Open Circuit Breaker for all external APIs & \texttt{TerrainSource} decouples input format from the analysis pipeline; adding GeoTIFF support requires only a new subclass with no changes to hydrology code. \\
    Error Handling \& Resilience & Every external call (OSM, Open-Meteo, NASA POWER, DB, Terrarium tiles) in try/except with fallback; warnings accumulate in \texttt{AnalysisResult.warnings[]} & The pipeline must return a usable result even when APIs are down. Machine-readable codes let the frontend show human-readable notices without coupling to error internals. \\
    Algorithms \& Complexity & D8 flow: $O(N)$ topological sort; Priority-Flood: $O(N \log N)$; IDW: $O(N \cdot k)$ per cell & D8 was chosen over D-infinity because integer single-direction pointers form a tree, enabling $O(N)$ BFS catchment delineation by reversing edges. \\
    Containerisation & \texttt{Dockerfile} + \texttt{docker-compose.yml} with healthcheck-gated DB startup & Reproducible environment for grading; compose also starts the optional PostgreSQL instance with correct credentials. \\
    Testing Strategy & Unit tests for all modules; integration tests for the full API pipeline; real KMZ fixtures & Unit tests validate IDW, SCS-CN, and IS~5477 formulae against hand-computed reference values; integration tests run the full pipeline on a known KMZ. \\
    Version Control / CI & GitHub Actions: black + ruff + pytest on every push; pre-commit hooks locally & Ruff caught an unused variable in \texttt{from\_tiles.py} at commit time; black enforces consistent style across the codebase automatically. \\
    \bottomrule
  \end{tabular}
\end{table*}

The three most significant CSD themes are \textbf{Error Handling and
Resilience}, \textbf{Algorithms and Complexity}, and \textbf{API Design}.

The fail-open resilience strategy was the most consequential engineering
decision. Because the system chains five external data sources, a
fail-closed approach would make it unavailable whenever any source had an
outage. By making each step independently fail-open and accumulating
\texttt{warnings[]} codes, the system always returns at least the
topographic analysis and candidate locations.

The choice of D8 over continuous flow models (D-infinity, MFD) was
deliberate. D8 assigns exactly one downstream neighbour per cell, making
the flow-direction grid a tree. This enables BFS catchment delineation in
$O(N)$ time by reversing tree edges --- no priority queue or
re-computation is required.

%% =====================================================================
\section{Results and Evaluation}
\label{sec:results}

The system was validated against a KMZ contour map of a test site in
Raipur district, Chhattisgarh (21.245°N, 81.301°E).

\begin{table}[h]
  \caption{Computed outputs for the Raipur test site}
  \label{tab:results}
  \begin{tabular}{p{3.8cm}p{3.4cm}}
    \toprule
    Metric & Value \\
    \midrule
    Catchment area & 19.5 ha \\
    Elevation range & 284--290 m \\
    Mean slope & 1.4° \\
    Annual rainfall (10-yr avg.) & 1,458 mm/yr \\
    Wet season (Jun--Sep) & 1,260 mm \\
    Curve Number (HSG-B) & CN 71 \\
    Annual runoff volume & 153,500 m³ \\
    Recommended water depth & 1.84 m \\
    Embankment height & 2.34 m \\
    Live storage & 61,400 m³ \\
    Total storage (IS~5477) & 90,000 m³ \\
    Design constrained by & Hydrology \\
    Topographic capacity & 224,200 m³ \\
    Wet-season fill & Yes (100\%) \\
    Candidates found & 10 \\
    \bottomrule
  \end{tabular}
\end{table}

The ``constrained by hydrology'' label indicates the terrain bowl can
store more than the catchment delivers; the pond is sized to capture
$\approx$40\% of annual runoff (IS~5477 planning value), leaving 60\%
for downstream flow and inter-annual variability.

\begin{figure}[h]
  \centering
  \includegraphics[width=0.5\linewidth]{Screenshot from 2026-09-27 22-38-44.png}
  \caption{Results overlay for the test site. Green polygon: delineated
    catchment; blue dots: candidate sites; red star: selected pour point.}
  \label{fig:results}
  \Description{Map showing the delineated catchment polygon and ten candidate
    pond sites for the Raipur test KMZ, alongside the results sidebar with
    the values from Table 4.}
\end{figure}

\subsection{Performance}

Response times on a single-core development machine (Intel Core i5, 8~GB RAM):

\begin{table}[h]
  \caption{Measured end-to-end response times}
  \label{tab:perf}
  \begin{tabular}{p{4.0cm}p{3.2cm}}
    \toprule
    Operation & Time \\
    \midrule
    KMZ parse + IDW DEM (19.5~ha) & 8--12 s \\
    Priority-Flood conditioning & 2--4 s \\
    D8 flow + candidates & 3--6 s \\
    OSM Overpass (network) & 0.5--2 s (cached: 0 ms) \\
    Open-Meteo rainfall & 0.1--0.5 s \\
    SCS-CN + IS~5477 & $<$0.1 s \\
    Watershed BFS + polygonise & 0.5--1 s \\
    \textbf{Total (KMZ run)} & \textbf{60--90 s} \\
    SRTM tile download (5 km²) & 1--3 s \\
    \textbf{Total (tile run)} & \textbf{30--60 s} \\
    \bottomrule
  \end{tabular}
\end{table}

The dominant cost is IDW interpolation on the 2~m grid ($O(N \cdot k)$,
$N \approx 5 \times 10^5$ cells for a 20~ha area). The SRTM tile path is
faster because the 38~m grid has 361$\times$ fewer cells for the same area.

%% =====================================================================
\section{Discussion and Limitations}
\label{sec:discussion}

\textbf{What worked well.} The fail-open pipeline architecture kept the
system functional throughout development even when OSM Overpass was
rate-limited and NASA POWER was temporarily unreachable. The D8 + BFS
approach delineated the correct catchment boundary on the test KMZ.
The SCS-CN method produced a plausible runoff depth of 895~mm against an
annual rainfall of 1,458~mm for HSG-B mixed-agriculture land, consistent
with published regional studies.

\textbf{SRTM resolution limitation.} The Draw Area feature uses SRTM
tiles at $\approx$38~m resolution. Small depressions (footprint
$<$0.15~ha) are invisible at this scale. In low-relief terrain ($<$2~m
elevation range), Priority-Flood may fill all depressions and return no
candidates. The KMZ upload path at 2~m resolution does not have this
issue. Users are informed via a warning banner and the help page.

\textbf{SCS-CN assumptions.} The Curve Number is applied uniformly using
a single HSG value (configurable, default HSG-B). Spatial soil
heterogeneity within the catchment is not modelled. The built-up fraction
adjustment uses the full DEM bounding box rather than the delineated
catchment (unavailable at that pipeline stage), introducing a slight
conservative bias.

\textbf{No government land verification.} No authoritative free API
exists for parcel-level land ownership data in India. OSM landuse
exclusion (buildings, residential/commercial) is a reasonable proxy but
will miss privately owned agricultural parcels.

\textbf{Single-threaded analysis.} Each request occupies one thread pool
worker for 60--90~s. A second concurrent analysis queues behind it. This
is acceptable for demonstration use but would require worker-pool scaling
or analysis result caching for a production deployment.

%% =====================================================================
\section{AI Tool Usage Declaration}
\label{sec:ai-usage}

AI assistance (Google Gemini and Claude Sonnet) was used throughout this
project for the following tasks:

\begin{itemize}
  \item \textbf{Architecture review}: LLMs reviewed High-Level Design
        documents and flagged inconsistencies (e.g., scoring formula
        mismatch, area-consistency threshold discrepancies).

  \item \textbf{Code generation}: Boilerplate for rainfall API clients,
        IS~5477 pond design, SQLModel DB models, and the Terrarium tile
        downloader was AI-generated. All generated code was read,
        understood, and manually verified against the relevant standards
        (SCS-CN USDA tables, IS~5477, IS~12169, OSM tagging conventions)
        before being committed.

  \item \textbf{Test generation}: Unit test skeletons were AI-generated;
        all expected numerical values in assertion statements were
        hand-computed or cross-referenced against published formula results.

  \item \textbf{Bug identification}: LLM review identified a critical bug
        where monthly rainfall was averaged to a daily value before
        applying the SCS-CN threshold, causing zero runoff output. The fix
        (applying SCS-CN directly to monthly totals) was verified against
        published worked examples.

  \item \textbf{Documentation}: This report template was reformatted and
        populated with AI assistance. All technical content reflects the
        actual system implementation and was verified by the author.
\end{itemize}

In all cases, no AI output was accepted without independent verification.
The author takes full responsibility for the correctness of all submitted
code, results, and statements in this report.

%% =====================================================================
\appendix

\section{Source Code, Demo, and Deployment}

\subsection*{Source Code (GitHub)}
The complete source code is publicly available at:
\begin{center}
  \url{https://github.com/harshitap1305/Village-Pond-Planning-System}
\end{center}

\subsection*{Live Deployed Application}
A live instance of the system is hosted at:
\begin{center}
  \url{https://YOUR-DEPLOYED-URL-HERE}  %% TODO: replace with your actual deployment URL
\end{center}
(e.g.\ a Render, Railway, or Hugging Face Spaces deployment)

\subsection*{Video Demo (YouTube)}
A walkthrough demonstration of the system is available at:
\begin{center}
  \url{https://www.youtube.com/watch?v=YOUR-VIDEO-ID}  %% TODO: replace after uploading
\end{center}
The video covers: uploading a KMZ file, using the Draw Area tool, reading
the pond design results, and the pour-point override feature.

\subsection*{Key Repository Directories}
\begin{itemize}
  \item \texttt{src/terrain/} --- KML/KMZ parsing and TerrainSource ABC
  \item \texttt{src/dem/} --- IDW DEM builder, Terrarium tile downloader,
        Priority-Flood conditioning, slope raster
  \item \texttt{src/hydrology/} --- D8 flow routing, BFS watershed, SCS-CN
        runoff estimation
  \item \texttt{src/catchment/} --- OSM exclusion masks, candidate scoring,
        IS~5477 pond design
  \item \texttt{src/external/} --- Rainfall API clients (Open-Meteo, NASA
        POWER) with retry and timeout
  \item \texttt{src/api/} --- FastAPI routes, analysis service orchestrator
  \item \texttt{frontend/} --- Leaflet SPA (\texttt{index.html}),
        documentation page (\texttt{help.html})
  \item \texttt{tests/unit/} --- Per-module unit tests with hand-computed
        reference values
  \item \texttt{tests/integration/} --- Full-pipeline API-level tests
  \item \texttt{docs/} --- Architecture document, module design docs,
        this report
\end{itemize}

\end{document}
\endinput
